\chapter{Communication Complexity — KL Divergence}
%%% generated by the blueprint pipeline from TCSlib.CommunicationComplexity.NewmanTheorem.KLDivergence
\section{Overview}

This module develops concrete sum formulas for the Kullback--Leibler divergence on finite
measurable spaces, bridging Mathlib's \texttt{InformationTheory.klDiv} with the PFR
real-valued \texttt{KLDiv} entropy API.  The main results give explicit pointwise
summation formulas for KL divergence between finite measures and PMFs, and show how they
specialize to the Boolean ($\mathrm{Bool}$) case.

%%%%%%%%%%%%%%%%%%%%%%%%%%%%%%%%%%%%%%%%%%%%%%%%%%%%%%%%%%%%%%%%%%%%%%%%%%%%%%%
\section{Declarations}

\begin{theorem}[KL divergence as a log-likelihood-ratio sum (absolutely continuous case)]
\lean{CommunicationComplexity.klDiv_eq_sum_llr_of_ac}
\label{CommunicationComplexity.klDiv_eq_sum_llr_of_ac}
\leanok
\uses{CommunicationComplexity.FiniteMeasureSpace}
\difficulty{3}
Let $\Omega$ be a finite measurable space, and let $\mu$ and $\nu$ be finite measures on
$\Omega$ with $\mu$ absolutely continuous with respect to $\nu$. Write
$\tfrac{d\mu}{d\nu}$ for the Radon–Nikodym derivative and let $\mathrm{llr}(\omega) =
\log \tfrac{d\mu}{d\nu}(\omega)$ be the log-likelihood ratio at $\omega$. Then the
Kullback–Leibler divergence satisfies
\[
  D_{\mathrm{KL}}(\mu \,\|\, \nu)
= \left[\; \sum_{\omega \in \Omega} \mu(\{\omega\})\,\mathrm{llr}(\omega) \;+\;
\nu(\Omega) \;-\; \mu(\Omega) \;\right]_{+},
\]
where $[\,x\,]_{+}$ denotes the value in $[0,\infty]$ obtained by viewing the real
number $x$ as a nonnegative extended real, negative values being sent to $0$.
\end{theorem}

\begin{theorem}[KL divergence on a finite space as a log-likelihood-ratio sum]
\lean{CommunicationComplexity.FiniteMeasureSpace.klDiv_eq_sum_llr}
\label{CommunicationComplexity.FiniteMeasureSpace.klDiv_eq_sum_llr}
\leanok
\uses{CommunicationComplexity.FiniteMeasureSpace, CommunicationComplexity.FiniteMeasureSpace.absolutelyContinuous_iff_forall_singletons, CommunicationComplexity.klDiv_eq_sum_llr_of_ac}
\difficulty{4}
Let $\Omega$ be a finite discrete measurable space, and let $\mu$ and $\nu$ be finite
measures on $\Omega$. If there is a point $\omega \in \Omega$ with $\nu(\{\omega\}) = 0$
but $\mu(\{\omega\}) \neq 0$, then the Kullback–Leibler divergence is infinite,
$D_{\mathrm{KL}}(\mu \,\|\, \nu) = \infty$. Otherwise,
\[
  D_{\mathrm{KL}}(\mu \,\|\, \nu)
= \left(\ \sum_{\omega \in \Omega} \mu(\{\omega\})\,\log\frac{d\mu}{d\nu}(\omega)\ +\
\nu(\Omega)\ -\ \mu(\Omega)\ \right)_{\!+},
\]
where $\log\frac{d\mu}{d\nu}$ denotes the log-likelihood ratio of $\mu$ with respect to
$\nu$, the measures are read as their real values, and $(x)_+$ takes the enclosed real
number as a nonnegative extended real, sending any negative value to $0$.
\end{theorem}

\begin{theorem}[KL divergence of probability mass functions]
\lean{CommunicationComplexity.FiniteMeasureSpace.pmf_klDiv_eq_sum_llr}
\label{CommunicationComplexity.FiniteMeasureSpace.pmf_klDiv_eq_sum_llr}
\leanok
\uses{CommunicationComplexity.FiniteMeasureSpace, CommunicationComplexity.FiniteMeasureSpace.klDiv_eq_sum_llr}
\difficulty{3}
Let $\Omega$ be a finite discrete measurable space, and let $p$ and $q$ be probability
mass functions on $\Omega$. Then the Kullback–Leibler divergence between the probability
measures they induce is
\[
  D_{\mathrm{KL}}(p \,\|\, q) =
  \begin{cases}
\infty & \text{if } q(\omega) = 0 \text{ and } p(\omega) \neq 0 \text{ for some } \omega
\in \Omega, \\[4pt]
\displaystyle\sum_{\omega \in \Omega} p(\omega)\,\log\!\frac{p(\omega)}{q(\omega)} &
\text{otherwise,}
  \end{cases}
\]
where in the second case the (nonnegative) real sum is read as an element of
$[0,\infty]$.
\end{theorem}

\begin{theorem}[Radon–Nikodym derivative as a ratio of point masses]
\lean{CommunicationComplexity.rnDeriv_toReal_eq_singleton_ratio}
\label{CommunicationComplexity.rnDeriv_toReal_eq_singleton_ratio}
\leanok
\uses{CommunicationComplexity.FiniteMeasureSpace}
\difficulty{4}
Let $\Omega$ be a finite measurable space carrying two finite measures $\mu$ and $\nu$
with $\mu \ll \nu$, and let $\omega \in \Omega$ be a point whose singleton has positive
mass, $\nu(\{\omega\}) \neq 0$. Then the Radon–Nikodym derivative of $\mu$ with respect
to $\nu$ takes at $\omega$ the value
\[
  \frac{d\mu}{d\nu}(\omega) = \frac{\mu(\{\omega\})}{\nu(\{\omega\})},
\]
where both sides are read as nonnegative real numbers.
\end{theorem}

\begin{theorem}[Singleton mass times log-likelihood ratio]
\lean{CommunicationComplexity.singleton_mass_mul_llr_eq_log_ratio}
\label{CommunicationComplexity.singleton_mass_mul_llr_eq_log_ratio}
\leanok
\uses{CommunicationComplexity.FiniteMeasureSpace, CommunicationComplexity.rnDeriv_toReal_eq_singleton_ratio}
\difficulty{3}
Let $\Omega$ be a finite discrete measurable space, let $\mu$ and $\nu$ be finite
measures on $\Omega$ with $\mu$ absolutely continuous with respect to $\nu$, and let
$\omega \in \Omega$. Writing $\mu(\{\omega\})$ and $\nu(\{\omega\})$ for the
(real-valued) masses of the singleton $\{\omega\}$ and $\log\frac{d\mu}{d\nu}$ for the
log-likelihood ratio of $\mu$ against $\nu$, one has
\[
  \mu(\{\omega\}) \cdot \log\frac{d\mu}{d\nu}(\omega)
  = \mu(\{\omega\}) \cdot \log\!\left(\frac{\mu(\{\omega\})}{\nu(\{\omega\})}\right).
\]
\end{theorem}

\begin{theorem}[KL divergence on a finite space as a sum of log-ratios]
\lean{CommunicationComplexity.FiniteMeasureSpace.klDiv_eq_sum_log}
\label{CommunicationComplexity.FiniteMeasureSpace.klDiv_eq_sum_log}
\leanok
\uses{CommunicationComplexity.FiniteMeasureSpace, CommunicationComplexity.FiniteMeasureSpace.absolutelyContinuous_iff_forall_singletons, CommunicationComplexity.FiniteMeasureSpace.klDiv_eq_sum_llr, CommunicationComplexity.singleton_mass_mul_llr_eq_log_ratio}
\difficulty{4}
Let $\Omega$ be a finite discrete measurable space, and let $\mu$ and $\nu$ be finite
measures on $\Omega$. If there is a point $\omega \in \Omega$ with $\nu(\{\omega\}) = 0$
but $\mu(\{\omega\}) \neq 0$, then the Kullback–Leibler divergence is infinite:
\[
  D_{\mathrm{KL}}(\mu \,\|\, \nu) = \infty.
\]
Otherwise,
\[
  D_{\mathrm{KL}}(\mu \,\|\, \nu)
  = \left(
      \sum_{\omega \in \Omega}
        \mu(\{\omega\}) \,\log\!\frac{\mu(\{\omega\})}{\nu(\{\omega\})}
      + \nu(\Omega) - \mu(\Omega)
    \right)_{+},
\]
where each $\mu(\{\omega\})$, $\nu(\{\omega\})$, $\mu(\Omega)$, $\nu(\Omega)$ denotes
the (real-valued) mass of the indicated set and $(x)_{+} = \max(x, 0)$ is the
nonnegative part.
\end{theorem}

\begin{theorem}[Finiteness of KL divergence against a full-support measure]
\lean{CommunicationComplexity.FiniteMeasureSpace.klDiv_ne_top_of_forall_toPMF_ne_zero}
\label{CommunicationComplexity.FiniteMeasureSpace.klDiv_ne_top_of_forall_toPMF_ne_zero}
\leanok
\uses{CommunicationComplexity.FiniteMeasureSpace, CommunicationComplexity.FiniteMeasureSpace.klDiv_eq_sum_log}
\difficulty{2}
Let $\Omega$ be a finite discrete measurable space, and let $\mu$ and $\nu$ be
probability measures on $\Omega$. If $\nu$ assigns positive mass to every point, meaning
$\nu(\{\omega\}) \neq 0$ for all $\omega \in \Omega$, then the Kullback–Leibler
divergence $D_{\mathrm{KL}}(\mu \,\|\, \nu)$ is finite.
\end{theorem}

\begin{theorem}[Kullback–Leibler divergence between probability mass functions]
\lean{CommunicationComplexity.FiniteMeasureSpace.pmf_klDiv_eq_sum_log}
\label{CommunicationComplexity.FiniteMeasureSpace.pmf_klDiv_eq_sum_log}
\leanok
\uses{CommunicationComplexity.FiniteMeasureSpace, CommunicationComplexity.FiniteMeasureSpace.klDiv_eq_sum_log}
\difficulty{3}
Let $\Omega$ be a finite discrete measurable space, and let $p$ and $q$ be probability
mass functions on $\Omega$, each regarded as the probability measure it induces. Then
the Kullback–Leibler divergence $D_{\mathrm{KL}}(p \,\|\, q)$, valued in $[0,\infty]$,
equals $+\infty$ whenever there exists some $\omega \in \Omega$ with $q(\omega) = 0$ and
$p(\omega) \neq 0$, and otherwise equals the value in $[0,\infty]$ obtained from
\[
  \sum_{\omega \in \Omega} p(\omega)\,\log\!\frac{p(\omega)}{q(\omega)},
\]
where each mass $p(\omega)$, $q(\omega)$ is taken as its real value and negative sums
are read as $0$.
\end{theorem}

\begin{theorem}[Real-valued and extended KL divergence agree on Bool]
\lean{ProbabilityTheory.toReal_klDiv_bool_eq_KLDiv}
\label{ProbabilityTheory.toReal_klDiv_bool_eq_KLDiv}
\leanok
\uses{CommunicationComplexity.FiniteMeasureSpace.klDiv_eq_sum_log}
\difficulty{4}
Let $\mu$ and $\nu$ be probability measures on the two-point space $\mathrm{Bool} =
\{\mathrm{false}, \mathrm{true}\}$, and suppose $\nu$ has full support, meaning
$\nu(\{b\}) > 0$ for each of the two points $b$. Writing $D_{\mathrm{KL}}(\mu \,\|\,
\nu) \in [0, \infty]$ for the Kullback–Leibler divergence valued in the extended
nonnegative reals, its real cast agrees with the real-valued Kullback–Leibler divergence
of the identity map, regarded as a random variable, under $\mu$ and under $\nu$; that
is,
\[
  \bigl(D_{\mathrm{KL}}(\mu \,\|\, \nu)\bigr)_{\bbr}
  = \mathrm{KL}[\,\mathrm{id};\mu \,\#\, \mathrm{id};\nu\,].
\]
\end{theorem}

\begin{theorem}[Real value of a conditioned Boolean Kullback–Leibler divergence]
\lean{ProbabilityTheory.toReal_klDiv_map_bool_eq_KLDiv_of_measureReal_ne_zero}
\label{ProbabilityTheory.toReal_klDiv_map_bool_eq_KLDiv_of_measureReal_ne_zero}
\leanok
\uses{ProbabilityTheory.toReal_klDiv_bool_eq_KLDiv}
\difficulty{3}
Let $\mu$ be a finite measure on a measurable space $\Omega$, let $X:\Omega\to\{0,1\}$
be a measurable Boolean random variable, and let $S\subseteq\Omega$ be a set whose
$\mu$-measure (as a real number) is nonzero, so that the conditional probability measure
$\mu(\cdot\mid S)$ is well defined. Let $\nu$ be a probability measure on $\{0,1\}$ that
assigns strictly positive mass to each of the two points. Then the extended-real
Kullback–Leibler divergence of the law $X_{\ast}\,\mu(\cdot\mid S)$ of $X$ under the
conditional measure relative to $\nu$, taken as a real number, coincides with its
explicit pointwise value:
\[
D_{\mathrm{KL}}\!\bigl(X_{\ast}\,\mu(\cdot\mid S)\ \big\|\ \nu\bigr)
= \sum_{b\in\{0,1\}} \bbP_{\mu(\cdot\mid S)}(X=b)\,
\log\frac{\bbP_{\mu(\cdot\mid S)}(X=b)}{\nu(\{b\})}.
\]
\end{theorem}
